\begin{pastexam}[label={px:ncu-108-dsa-06}]{中央資工所 108 學年度 第 6 題}
（複選；每題 5 分，答錯倒扣 1 分）

Consider an empty hash table with 10 buckets and each bucket has 2 slots. Suppose that the hash
function is h(n) = n \% 10, and the following numbers, 6 \ 16 \ 22 \ 45 \ 54 \ 25 \ 7 \ 75 \ 5 \ 108,
are sequentially inserted into the hash table. Which of the following statements are true?
\begin{enumerate}[label=(\Alph*)]
\item if linear probing is adopted to handle overflow, the summation of the numbers in the full
buckets of the hash table is 287.
\item if linear probing is adopted to handle overflow, the summation of the numbers in the full
buckets of the hash table is 172.
\item if quadratic probing is adopted to handle overflow, the summation of the numbers in the full
buckets of the hash table is 221.
\item Compared with quadratic probing, linear probing requires fewer bucket accesses (in average).
\item None of the above.
\end{enumerate}
\end{pastexam}

\begin{pastsoln}
答案：(A)(C)。

每個桶有兩格，裝得下兩個鍵；「full bucket」是兩格都滿的桶。 探測以桶為單位：$h(n)$ 號桶滿了才照探測
序列試下一個桶。 雜湊值就是個位數：
\begin{center}
\small
\begin{tabular}{c*{10}{c}}
\toprule
$n$ & 6 & 16 & 22 & 45 & 54 & 25 & 7 & 75 & 5 & 108 \\
$h(n)$ & 6 & 6 & 2 & 5 & 4 & 5 & 7 & 5 & 5 & 8 \\
\bottomrule
\end{tabular}
\end{center}
前七個數都直接放進自己的桶：$6$ 號桶放 $6, 16$（滿），$5$ 號桶放 $45, 25$（滿），$2, 4, 7$ 號桶各一個。
差別在後面三個。

\textbf{線性探測}（依序試 $h, h+1, h+2, \ldots$）：
\begin{itemize}
\item $75$：$5$ 號滿、$6$ 號滿，放進 $7$ 號（和 $7$ 一起，滿）。
\item $5$：$5, 6, 7$ 號都滿，放進 $8$ 號。
\item $108$：$8$ 號還有一格，放進去（滿）。
\end{itemize}
滿的桶是 $5$ 號 $\{45, 25\}$、$6$ 號 $\{6, 16\}$、$7$ 號 $\{7, 75\}$、$8$ 號 $\{5, 108\}$，總和
$70 + 22 + 82 + 113 = 287$。 (A) 對，(B) 錯。

\textbf{二次探測}。 資結書第~12.3~節的二次探測是 $h + i^2$（試 $h, h+1, h+4, h+9, \ldots$）；另一種常見的
寫法是正負交替 $h, h+1^2, h-1^2, h+2^2, h-2^2, \ldots$（都取 $\bmod 10$）。 兩種的結果不同：
\begin{center}
\small
\begin{tabular}{lll}
\toprule
 & $h + i^2$ & $h \pm i^2$ \\
\midrule
$75$（$h = 5$） & $5$ 滿、$6$ 滿、$9$：放進 $9$ 號 & $5$ 滿、$6$ 滿、$4$：放進 $4$ 號（和 $54$，滿） \\
$5$（$h = 5$） & $5$ 滿、$6$ 滿、$9$：放進 $9$ 號（滿） & $5$、$6$、$4$ 都滿，$5 + 4 = 9$：放進 $9$ 號 \\
$108$（$h = 8$） & 放進 $8$ 號 & 放進 $8$ 號 \\
\midrule
滿的桶 & $5, 6, 9$ 號 & $4, 5, 6$ 號 \\
總和 & $70 + 22 + 80 = 172$ & $129 + 70 + 22 = 221$ \\
\bottomrule
\end{tabular}
\end{center}
（正負交替時先試 $h - 1^2$ 也一樣：$75$ 進 $4$ 號，$5$ 改進 $5 - 4 = 1$ 號，滿的桶仍是 $4, 5, 6$ 號，$221$。）
選項 (C) 的 $221$ 只有正負交替的版本算得出來，出題者用的是這個版本，(C) 對。 若照 $h + i^2$ 算，總和
是 $172$，(C) 就是錯的；作答時最好寫出採用的探測序列。

\textbf{(D)}。 錯。 線性探測有主要群聚：連成一段的滿格越長，越容易被接上而變得更長，平均要看的桶比較多
（資結書第~12.3、12.4~節；失敗搜尋約 $\frac12\bigl(1 + 1/(1-\alpha)^2\bigr)$ 次，比均勻雜湊的
$1/(1 - \alpha)$ 多）。 二次探測的跳距越來越大，沒有主要群聚，平均存取次數比線性探測少，不是多。

\textbf{(E)} (A)、(C) 對，所以錯。
\end{pastsoln}
